\documentclass[10pt,a4paper]{amsart}
\usepackage[margin=19mm]{geometry}
\usepackage{amsmath,amssymb,mathtools}
\usepackage[scr=rsfs,cal=euler]{mathalfa}
\usepackage{xfrac}
\usepackage[unicode,hidelinks,bookmarks=false]{hyperref}
\newcommand{\ie}{i.e.\ }
\newcommand{\cf}{cf.\ }
\newcommand{\resp}{respectively\ }
\newcommand{\Subsection}{\subsection}
\providecommand{\nobreakdash}{\nobreak\hbox{-}}
\setlength{\emergencystretch}{2em}
\multlinegap0pt
\numberwithin{equation}{section}
\newtheorem{theo}{Theorem}[section]
\newtheorem{theorem}[theo]{Theorem}
\newtheorem{lemma}[theo]{Lemma}
\newtheorem{definition}[theo]{Definition}
\newtheorem{proposition}[theo]{Proposition}
\newtheorem{remark}[theo]{Remark}
\newtheorem{corollary}[theo]{Corollary}
\newtheorem{example}[theo]{Example}
\allowdisplaybreaks
\title[Multiple positive normalized solutions for Kirchhoff type sy]{Multiple positive normalized solutions for Kirchhoff type system with van der Waals type potentials: Theorem1.2}
\author{Zhewen Chen; Muzi Li}
\date{Source: Electronic Journal of Differential Equations 2026, No. 06, pp. 1--20; DOI 10.58997/ejde.2026.06; publisher version}
\begin{document}
\maketitle
\noindent\emph{Electronic Journal of Differential Equations}, Vol. 2026 (2026), No. 06, pp. 1--20.\newline
ISSN: 1072-6691. DOI \href{https://doi.org/10.58997/ejde.2026.06}{10.58997/ejde.2026.06}. This work is licensed under a CC BY 4.0 license.\par\smallskip
\noindent\textit{Editorial scope.} Counted claim: original Theorem 1.2 (the first main theorem of the article; source label thm1.1, source lines 345-350) together with its complete original proof at source lines 844-892, delivered as the single primary proof body. Necessary complete context is retained uncounted: the whole of Section 1 (the system, the mass constraint, the Riesz potentials, Definition 1.1 of a normalized ground state, assumptions (A1) and (A2), the Pohozaev manifold and the level sigma(d\_1,d\_2)) and the whole of Section 2 (Lemmas 2.1-2.12 with their complete proofs), together with Lemma 3.1 and its complete proof at source lines 815-842. Theorems 1.3 and 1.4 and their Sections 4 and 5 are not selected and are not used to inflate H. Slips kept literal and recorded without mathematical repair: (i) the lemma header at source line 568 contains a malformed citation command whose opening brace was printed as the letter e; because that byte sequence cannot be delivered as a well-formed citation and would also make the compiler report an undefined citation, the single header line is supplied editorially with exactly the published rendering, namely the lemma with the bracket note containing the printed question mark and the stray word linjie, while the lemma statement body at source lines 569-578 is copied byte-exact and its label is kept, so that Lemma 2.12 still refers to it and it prints its published number 2.7; (ii) the reference to the undefined name Thm1.3 in Remark 1.5, printed as a double question mark in the published PDF, is delivered as a hyperlink whose visible text is that same literal printed slip; (iii) the references in Remark 1.6 to the Lemmas numbered 4.1 and 4.2, whose statements and proofs lie outside the counted unit in Section 4, are delivered as hyperlinks carrying their published numbers 4.1 and 4.2. No mathematics is written, repaired or re-derived; all mathematics is native text with no image dependency.
\section{Introduction and main results}

In this article, we study the existence of the positive normalized solutions for
Kirchhoff type system with van der Waals type potentials
\begin{equation}\label{yy}
\begin{gathered}
-(a+b\int_{\mathbb{R}^N}|\nabla u_1|^2dx)\Delta u_1
=\lambda_1 u_1+\mu_1(I_\alpha\ast|u_1|^{p_1})|u_1|^{p_1-2}u_1+
 \Theta r_1(I_\beta\ast|u_2|^{r_2})|u_1|^{r_1-2}u_1, \\
-(a+b\int_{\mathbb{R}^N}|\nabla u_2|^2dx)\Delta u_2
=\lambda_2 u_2+\mu_2(I_\alpha\ast|u_2|^{p_2})|u_2|^{p_2-2}u_2+
 \Theta r_2(I_\beta\ast|u_1|^{r_1})|u_2|^{r_2-2}u_2,
\end{gathered}
\end{equation}
with the $L^2$-mass constraint
\begin{equation}\label{yues}
 \int_{\mathbb{R}^N}|u_1|^2dx=d_1>0,\quad
 \int_{\mathbb{R}^N}|u_2|^2dx=d_2>0,
\end{equation}
where $N=3,4$, $a,b>0$, $\alpha,\beta\in(0,N)$, $I_\alpha,I_\beta$ are the Riesz
potentials defined for every $x\in\mathbb{R}^N\backslash\{0\}$ by
\begin{equation}\label{gamma}
 I_\alpha(x):=\frac{A_\alpha(N)}{|x|^{N-\alpha}},\quad
 A_\alpha(N):=\frac{\Gamma(\frac{N-\alpha}{2})}{\Gamma(\frac{\alpha}{2})\pi^{N/2}
 2^\alpha}
\end{equation}
with $\Gamma$ denoting the Gamma function. Throughout this paper, we always require
that $a,b,\mu_1,\mu_2,\Theta>0$, and assume that $2<p_1,p_2
<2+\frac{4}{N}$, $r_1,r_2>\frac{N+\beta}{N}$,
$2\cdot\frac{N+\beta}{N}<r_1+r_2<2\cdot2_\beta^*$.

Because of the appearance of the term $\int_{\mathbb{R}^N}|\nabla u|^2dx$,
\eqref{yy} is regard as a nonlocal problem, which implies that equation \eqref{yy}
is not a pointwise identity. Moreover, this phenomenon also leads to some mathematical
difficulties that make the study of \eqref{yy} more interesting.
Problem \eqref{yy} originates
from the stationary analog of the equation
\[
\rho\frac{\partial^2 u}{\partial t^2}
-\Big(\frac{Y_0}{k}+\frac{R}{2J}\int_{0}^{Q}|\frac{\partial u}{\partial x}|^2dx\Big)
\frac{\partial^2 u}{\partial x^2}=0,
\]
which was proposed by Kirchhoff \cite{Kf4} in 1883, and being an extension of the
classical D'Alembert's wave equations for free vibration of elastic strings.
Kirchhoff's model takes into account the changes in length of the string produced by
transverse vibrations. In the last decades, Because of the strong background in physics,
starting with the framework of Lions \cite{Li4}, many mathematicians have established
many interesting conclusions about Kirchhoff-type problems
\cite{Kf1,Kf2,Kf3,Kf4,Kf5,Kf6,ygg,guo1,guo2}.

In \eqref{yy}, if $\lambda\in\mathbb{R}$ is fixed, then we call \eqref{yy} the fixed
frequency problem. One can adopt the traditional variational method,
looking for critical points of $F_\lambda(u_1,u_2)$, or fixed point theory,
bifurcation, topological methods, Nehari manifold method and Lyapunov-Schmidt reduction,
where
\begin{align*}
F_\lambda(u_1,u_2)
&:=\frac{a}{2}\sum_{i=1}^2\| \nabla u_i\|_2^2
 +\frac{b}{4}\sum_{i=1}^2\| \nabla u_i\|_2^4-\sum_{i=1}^2\lambda_i\| u_i\|_2^2
 -\sum_{i=1}^2\frac{\mu_i}{2p_i}\int_{\mathbb{R}^N}
 (I_\alpha\ast|u_i|^{p_i})|u_i|^{p_i}dx\\
&\quad -\Theta\int_{\mathbb{R}^N}(I_\beta\ast|u_1|^{p_1})|u_2|^{p_2}dx.
\end{align*}
In recent decades, because of the application to physics, mathematicians are
interested in solutions that satisfy $L^2-$mass constraint \eqref{yues}.
In this direction, the mass $d_1,d_2>0$ is prescribed, the frequency $\lambda_i$
cannot be determined a priori, but is a part of unknown which appears as Lagrange
multipliers.
In this case, mathematicians often call \eqref{yy}-\eqref{yues} the fixed mass
problem and the solution is called a normalized solution.
One can get a normalized solution to problem \eqref{yy} by looking for a critical
point of the  functional
\begin{align*}
&F(u_1,u_2) \\
&:=\frac{a}{2}\sum_{i=1}^2\| \nabla u_i\|_2^2
 +\frac{b}{4}\sum_{i=1}^2\| \nabla u_i\|_2^4-
\sum_{i=1}^2\frac{\mu_i}{2p_i}\int_{\mathbb{R}^N}(I_\alpha\ast|u_i|^{p_i})|u_i|^{p_i}dx
-\Theta\int_{\mathbb{R}^N}(I_\beta\ast|u_1|^{p_1})|u_2|^{p_2}dx
\end{align*}
constrained on $S(d_1)\times S(d_2)$, where
$S(d):=\{u\in H^1(\mathbb{R}^N): \int_{\mathbb{R}^N}|u|^2dx=d>0\}$.
It is standard to check that $F\in C^1$.

For the fixed mass problem, the $L^2$-mass constraint \eqref{yues} presents some
mathematical difficulties. As opposed to the fixed
frequency problem, the fixed mass problem will have many technical difficulties when
dealing with it in a variational framework:
(I) the Nehari manifold method is inapplicable;
(II) the Lagrange multipliers must be controlled;
(III) for the fixed frequency problem, usually a
nontrivial weak limit is also a solution.
However, for the fixed mass problem, even though the weak limit is nontrivial, the
$L^2$-mass constraint may be not satisfied;
(IV) the $L^2$-mass critical exponent seriously affects the geometric structure of the
functional.

The equation
\begin{equation}\label{yyu5}
 i\partial_t\varPsi+\Delta \varPsi+(V\ast|\varPsi|^p)|\varPsi|^{p-2}\varPsi=0,    \text{in} \mathbb{R}^+\times\mathbb{R}^N,\\
 \end{equation}
has $\varPsi:\mathbb{R}^+\times\mathbb{R}^N\to \mathbb{C}$ is a complex
valued function, $V(x)=\delta_{\zeta}\frac{1}{|x|^\zeta}\pm
 \delta_{\theta}\frac{1}{|x|^\theta}$($\alpha,\beta>1$)
is the van der Waals type potential (see \cite{VWp3,BZd27}).
The van der Waals coefficients
$\delta_6$, $\delta_8$ and $\delta_{10}$ of alkaline-earth interactions
calculated by Porsev and Derevianko using relativistic many-body perturbation
theory are believed to be accurate to $1/100$ (see\cite{cfrv}).
As one of the van der Waals type potentials the Lennard-Jones potential
\[
 V_{LJ}(\nu)=\pm\xi(\frac{1}{\nu^\zeta }-\frac{1}{\nu^\theta})
\]
with $\zeta =6$ and $\theta=12$, is often used as an approximate model for the
isotropic part of a total van der
 Waals force as a function of distance (see \cite{VWp3,BZd27}).

When $a=1$, $b=0$ and the response function is a delta function, i.e.\
$I_\alpha(x)=I_\beta(x)=\delta(x)$, the nonlinear response is local and
problem \eqref{yy} with
prescribed mass turns out to be
\begin{equation}\label{hend}
\begin{gathered}
-\Delta u_1=\lambda_1 u_1+\mu_1|u_1|^{p_1-2}u_1+
 \Theta r_1|u_2|^{r_2}|u_1|^{r_1-2}u_1, \\
-\Delta u_2=\lambda_2 u_2+\mu_2|u_2|^{p_2-2}u_2+
 \Theta r_2|u_1|^{r_1}|u_2|^{r_2-2}u_2, \\
\end{gathered}
\end{equation}
with the $L^2$-mass constraint
$\int_{\mathbb{R}^N}|u_1|^2dx=d_1>0$, $\int_{\mathbb{R}^N}|u_2|^2dx=d_2>0$.
In recent years, mathematicians have drawn rich conclusions
concerning the existence, multiplicity and qualitative properties of
the normalized solutions of system \eqref{hend}
(see \cite{Huij1,Huij2,Huij3,Huij4,Huij5,Huij6}).
Gou et al.\ \cite{Huij6} showed that
the system admits two normalized solutions under the conditions that
$2<p_1,p_2<2+\frac{4}{N}<r_1+r_2<2^*=\frac{2N}{N-2}$
(resp.\ $2+\frac{4}{N}<r_1+r_2<2+\frac{4}{N}<p_1,p_2<2^*$)
by using minimizing methods, Poho\v{z}aev type manifold, Schwarz rearrangements
and mountain pass theorem.
Bartsch et al.\ \cite{Huij4} focused on the choice $p_1=p_2=2^*$, and allowing
$r_1+r_2$ to be mass-subcritical, mass-critical or mass-supercritical.
The authors proved the existence and non-existence of normalized ground state for
different ranges of $\Theta$.

When $a,b>0$ and the response function is a delta function,
i.e.\ $I_\alpha(x)=I_\beta(x)=\delta(x)$,
system \eqref{yy} becomes
\begin{equation}\label{sss8}
\begin{gathered}
 -(a+b\int_{\mathbb{R}^N}|\nabla u_1|^2dx)\Delta u_1=\lambda_1 u_1+\mu_1|u_1|^{p_1-2}u_1+
 \Theta r_1|u_2|^{r_2}|u_1|^{r_1-2}u_1   x\in\mathbb{R}^N,\\
-(a+b\int_{\mathbb{R}^N}|\nabla u_2|^2dx)\Delta u_2=\lambda_2 u_2+\mu_2|u_2|^{p_2-2}u_2+
 \Theta r_2|u_1|^{r_1}|u_2|^{r_2-2}u_2   x\in\mathbb{R}^N.\\
 \end{gathered}
\end{equation}
When $N\leq3$ and $2<p_1,p_2,r_1+r_2<2+\frac{8}{N}$, Cao et al.\
 \cite{kf2} proved that system \eqref{sss8} admits a positive
normalized solution. When $N=2,3$, Yang \cite{Kf3} showed that system
\eqref{sss8} admits a normalized ground state under the condition that
$2+\frac{8}{N}<p_1,p_2,r_1+r_2<2^*$ and $2<r_1+r_2<2+\frac{8}{N}<p_1,p_2<2^*$,
respectively.

When $a=1$, $b=0$ and the response function is a delta function, i.e.\
$I_\alpha(x)=I_\beta(x)=\delta(x)$,
system \eqref{yy} becomes the   Choquard system (Hartree system)
\begin{equation}\label{ChqH}
 \begin{gathered}
-\Delta u_1=\lambda_1 u_1+\mu_1(I_\alpha\ast|u_1|^{p_1})|u_1|^{p_1-2}u_1
+\Theta r_1(I_\alpha\ast|u_2|^{r_2})|u_1|^{r_1-2}u_1   x\in\mathbb{R}^N,\\
-\Delta u_2=\lambda_2 u_2+\mu_2(I_\alpha\ast|u_2|^{p_2})|u_2|^{p_2-2}u_2
+\Theta r_2(I_\alpha\ast|u_1|^{r_1})|u_2|^{r_2-2}u_2   x\in\mathbb{R}^N.
 \end{gathered}
\end{equation}
When $\mu_1,\mu_2,\Theta>0$ and $\frac{N+\alpha}{N}<r_1,r_2$, Geng et al.\
\cite{JFx} proved that system \eqref{ChqH}
has a normalized ground state under the condition that $N\geq3$ and
$\frac{N+\alpha}{N}<p_1,p_2<\frac{N+\alpha+2}{N}<r_1,r_2<2_\alpha^*$
by using Schwartz rearrangement. In addition, they proved that the system
\eqref{ChqH} has a second solution
for $N=3$ and $p_1=p_2=\alpha=2$.
They also proved that  \eqref{ChqH} has a second solution
for $N=5$, $\frac{N+\alpha}{N}<r_1,r_2<\frac{N+\alpha+2}{N}<p_1,p_2<2_\alpha^*$ and
$p_1=p_2=\alpha=2$. For
Hardy-Littlewood-Sobolev critical case, i.e.\ $p_1=p_2=2_\alpha^*$, Zhang et al.\
\cite{linjie} proved that system \eqref{ChqH}
has a normalized ground state for different ranges of $\Theta$ when $r_1+r_2$ is set
to be mass subcritical, mass critical and mass supercritical, respectively.

However, as far as we know, for the case of $a,b>0$ and the response functions are
different Riesz potential functions,
the existence of the solution of the system is still unknown.

\begin{definition}\label{df82} \rm
 $(u_1,u_2)$ is a \emph{normalized ground state} to \eqref{yy}-\eqref{yues}
if $F'|_{S(d_1)\times S(d_2)}(u_1,u_2)=0$ and
\[
 F(u_1,u_2)=\inf\{F(v_1,v_2): (v_1,v_2)\in S(d_1)\times S(d_2),
  F'|_{S(d_1)\times S(d_2)}(v_1,v_2)=0\}.
\]
Furthermore, $(w_1,w_2)$ is a \emph{high energy normalized solution}
to \eqref{yy}-\eqref{yues} if $F'|_{S(d_1)\times S(d_2)}(w_1,w_2)=0$ and
\[
 F(w_1,w_2)>\inf\{F(v_1,v_2): (v_1,v_2)\in S(d_1)\times S(d_2),
 F'|_{S(d_1)\times S(d_2)}(v_1,v_2)=0\}.
\]
\end{definition}

In this article, we demonstrate our main results in the following two cases:
\begin{itemize}
\item[(A1)] $\frac{N+\alpha}{N}<p_1,p_2<\frac{N+\alpha+2}{N}$,
$r_1,r_2>\frac{N+\beta}{N}$ and
$2\cdot\frac{N+\beta}{N}<r_1+r_2<2\cdot\frac{N+\beta+4}{N}$.

\item[(A2)]  $\frac{N+\alpha}{N}<p_1,p_2<\frac{N+\alpha+2}{N}$,
$r_1,r_2>\frac{N+\beta}{N}$ and $2\cdot\frac{N+\beta+4}{N}<r_1+r_2<2\cdot2_\beta^*$.
\end{itemize}
To restore some compactness, we search for critical points of $F$ constrained on
$S_r(d_1)\times S_r(d_2)$,
where $S_r(d):=\{u\in H_r^1(\mathbb{R}^N):\| u\|_2^2=d\}$ and
\[
H_r^1(\mathbb{R}^N)=\{u\in H^1(\mathbb{R}^N):
 u \text{is a radially symmetric function}.
\] 
For case of ($A_2$), $F$ is unbounded below in $S_r(d_1)\times S_r(d_2)$.
Hence, we need the Poho\v{z}aev manifold
$\mathcal{P}(d_1,d_2):=\{(u_1,u_2)\in S_r(d_1)\times S_r(d_2): P(u_1,u_2)=0\}$, where
\begin{align*}
P(u_1,u_2)
&=a\sum_{i=1}^2\| \nabla u_i\|_2^2+b\sum_{i=1}^2\| \nabla u_i\|_2^4
-\sum_{i=1}^2\frac{\mu_i}{2p_i}(Np_i-N-\alpha)
 \int_{\mathbb{R}^N}(I_\alpha\ast|u_i|^{p_i})|u_i|^{p_i}dx\\
&\quad-\Theta\frac{Nr-2(N+\beta)}{2}\int_{\mathbb{R}^N}(I_\beta\ast|u_1|^{r_1})
|u_2|^{r_2}dx
\end{align*}
with $r=r_1+r_2$. We define $W(k):=\{(u_1,u_2)\in S_r(d_1)\times
S_r(d_2):\| \nabla u_1\|_2^2+\| \nabla u_2\|_2^2<k\}$
and
\begin{equation}\label{jxh25}
 \sigma(d_1,d_2):=\inf_{W(k)}F(u_1,u_2)<0.
\end{equation}
For each $u\in S_r(d)$, we define
\[
 \tau\star u_i:=e^{\frac{N}{2}\tau}u_i(e^{\tau}x), \quad i=1,2.
\]
Hence,
\begin{equation}
\begin{aligned}
\Psi_{u_1,u_2}(\tau)
&:=F(\tau\star u_1,\tau\star u_2)\\
&=\frac{a}{2}e^{2\tau}\sum_{i=1}^2\|\nabla u_i\|_2^2
 +\frac{b}{4}e^{4\tau}\sum_{i=1}^2\|\nabla u_i\|_2^4\\
&\quad-\sum_{i=1}^2\frac{\mu_i}{2p_i}e^{(Np_i-N-\alpha)\tau}
 \int_{\mathbb{R}^N}(I_\alpha\ast|u_i|^{p_i})|u_i|^{p_i}dx\\
&\quad -\Theta e^{\frac{Nr}{2}-N-\beta}
 \int_{\mathbb{R}^N}(I_\beta\ast|u_1|^{r_1})|u_2|^{r_2}dx.
\end{aligned}
\end{equation}
By direct computations, we have $(\Psi_{u_1,u_2})'(0)=P(u_1,u_2)$.
We divide $\mathcal{P}(d_1,d_2)$ into 3 parts,
\[
\mathcal{P}(d_1,d_2)= \mathcal{P}_{-}(d_1,d_2)\cup\mathcal{P}_{0}(d_1,d_2)
\cup\mathcal{P}_{+}(d_1,d_2),
\]
 where
\begin{gather*}
\mathcal{P}_{-}(d_1,d_2)=\{(u_1,u_2)\in\mathcal{P}(d_1,d_2):(\Psi_{u_1,u_2})''(0)<0\},\\
\mathcal{P}_{0}(d_1,d_2)=\{(u_1,u_2)\in\mathcal{P}(d_1,d_2):(\Psi_{u_1,u_2})''(0)=0\},\\
\mathcal{P}_{+}(d_1,d_2)=\{(u_1,u_2)\in\mathcal{P}(d_1,d_2):(\Psi_{u_1,u_2})''(0)>0\},
\end{gather*}
and
\begin{align*}
(\Psi_{u_1,u_2})''(0)
&=2a\sum_{i=1}^2\|\nabla u_i\|_2^2+4b\sum_{i=1}^2\|\nabla u_i\|_2^4 \\
&\quad -\sum_{i=1}^2\frac{\mu_i}{2p_i}(Np_i-N-\alpha)^2
\int_{\mathbb{R}^N}(I_\alpha\ast|u_i|^{p_i})|u_i|^{p_i}dx \\
&\quad -\Theta(\frac{Nr}{2}-N-\beta)^2
 \int_{\mathbb{R}^N}(I_\beta\ast|u_1|^{r_1})|u_2|^{r_2}dx.
\end{align*}


\setcounter{theo}{1}
\begin{theorem} \label{thm1.1}
Assume that {\rm (A1)} holds. Then, every minimizing sequence of \eqref{jxh25}
is compact, up to translation, in
$H^1_r(\mathbb{R}^N)\times H^1_r(\mathbb{R}^N)$.
Moreover, system \eqref{yy} has a positive normalized ground state.
\end{theorem}


\begin{theorem} \label{thm1.2}
Assume that {\rm (A2)} holds and $0<\alpha-2\leq\beta<\alpha<N$.
Then, there exist $k_0=k_0(d_1,d_2)>0$, $\Theta_*=\Theta_*(d_1,d_2)>0$, such that
for any $0<\Theta\leq\Theta_*$,
system \eqref{yy}-\eqref{yues} has a positive radial solution
$(\upsilon_1,\upsilon_2)\in W(k_0)$ at negative
level $F(\upsilon_1,\upsilon_2)<0$ for some $\lambda_1,\lambda_2<0$.
 \end{theorem}

\begin{theorem} \label{thm1.3}
Assume that {\rm (A2)} holds,and $a=1$, $b=0$, $N=3$ and
$p_1=p_2=\alpha=\beta=r_1=r_2=2$. Then, there exist
$k_0= k_0(d_1,d_2)>0$, $\Theta_*=\Theta_*(d_1,d_2)>0$, such that for any
$0<\Theta\leq\Theta_*$,
system \eqref{yy}-\eqref{yues} has a second positive radial solution
$(u_1,u_2)\in W(k_0)$ at positive level $F(u_1,u_2)>0$
 for some $\lambda_1,\lambda_2<0$.
\end{theorem}

% page 4

 \begin{remark}  \label{rmk1.1} \rm
Definition \ref{df82}, Theorem \ref{thm1.2} and \href{https://ejde.math.txstate.edu/Volumes/2026/06/chen.pdf}{??}
indicate that \eqref{yy}-\eqref{yues} admit a normalized ground state at negative
level $F(\upsilon_1,\upsilon_2)<0$ and a high energy normalized solution at positive 
level $F(u_1,u_2)>0$, under the condition that
(A2), $a=1$, $b=0$, $N=3$ and $p_1=p_2=\alpha=\beta=r_1=r_2=2$ hold.
\end{remark}

\begin{remark} \label{rmk1.2} \rm
For nonlinear classical Choquard system \ref{ChqH}, the mass critical exponent is 
$\frac{N+\alpha+2}{N}$. In Kirchhoff-Choquard system,
by the appearance of the $\int_{\mathbb{R}^N}|\nabla u|^4dx$ term, the mass critical 
exponent becomes $\frac{N+\alpha+4}{N}$. In this
paper, if $\frac{N+\alpha}{N}<p_1,p_2<\frac{N+\alpha+4}{N}$, then Lemma 
\href{https://ejde.math.txstate.edu/Volumes/2026/06/chen.pdf}{4.1} and Lemma \href{https://ejde.math.txstate.edu/Volumes/2026/06/chen.pdf}{4.2} do not hold.
\end{remark}

\begin{remark} \label{rmk1.3} \rm
Since we do not know whether the inequality
 \[
\| \nabla \{u_1,u_2\}^*\|_2^4\leq\| \nabla u_1\|_2^4+\| \nabla u_2\|_2^4
 \]
holds, we can not use the Schwartz rearrangement method to prove the compactness
of minimized sequence in this paper.
\end{remark}


From a physical point of view, it is of great importance to study the solution of 
 problem  \eqref{yy}, as confirmed in  \cite{VWp3,BZd27,cfrv}.
 We emphasize that this study seems to be the first contribution regarding existence 
of normalized ground states for a Kirchhoff system with van der Waals type potentials.

The rest of this article organized as follows: 
In Section 2, we present some preliminaries. 
In Section 3, we prove the existence
 of normalized ground states under the purely mass subcritical case. 
In Section 4, we prove the existence of the
 first solution, which is a local minimizer. 
In Section 5, the second solution is proved by using mountain pass theorem.
\smallskip

\noindent\textbf{Notation:} 
 $L^s(\mathbb{R}^N)$ is the Lebesgue space with the norms 
 $\| u\| _s=(\int_{\mathbb{R}^N}|u|^sdx)^{1/s}$,  $1<s<\infty$.\\
$H^1(\mathbb{R}^N)$ is the usual Sobolev space with norm 
$\| u\|_{H^1(\mathbb{R}^N)}=
 (\int_{\mathbb{R}^N}|\nabla u|^2+|u|^2dx)^{1/2}$.



\section{Preliminaries}

\begin{lemma}\label{L22}[Hardy-Littlewood-Sobolev inequality \cite{lieb}]
Let $N\geq1$, $p, r>1$, and $0<\beta<N$ with $1/p+(N-\beta)/N+1/r=2$. 
Let $u\in L^p(\mathbb{R}^N)$ and $v\in L^r(\mathbb{R}^N)$. 
Then there exists a sharp constant $C(N,p,\beta)$, independent of $u$ and $v$, 
such that
\[
\big|\int_{\mathbb{R}^N}\int_{\mathbb{R}^N}\frac{u(x)v(y)}{|x-y|^{N-\beta}}dxdy\big|
\leq C_{N,p,r,\beta}\| u\|_p \| v\| _r.
\]
If $p=r=\frac{2N}{N+\beta}$, then
\[
 C_{N,p,r,\beta}=C_{N,\beta}
 =\pi^{\frac{N-\beta}{2}}
 \frac{\Gamma(\frac{\beta}{2})}{\Gamma(\frac{N+\beta}{2})}
 \Big\{\frac{\Gamma(\frac{N}{2})}{\Gamma(N)}\Big\}^{-\frac{\beta}{N}}.
\]
\end{lemma}

\begin{lemma}[Gagliardo-Nirenberg inequality of Power type \cite{GN23}] \label{Gn75}
 Let $N\geq1$ and $2\leq p<2^*$, then the following sharp Gagliardo-Nirenberg 
inequality
 \begin{equation}\label{GN}
\| u\|_p\leq C_{N,p}\| u\|_2^{1-\delta_p}\| \nabla u\|_2^{\delta_p}
\end{equation}
holds for any $u\in H^1(\mathbb{R}^N)$, where $\delta_p=\frac{N}{2}-\frac{N}{p}$, 
the sharp constant $C_{N,p}$ is
\[
 C_{N,p}^p=\frac{2p}{2N+(2-N)p}(\frac{2N+(2-N)p}{N(p-2)})^{\frac{N(p-2)}{4}}
 \frac{1}{\|Q_p\|_2^{p-2}}
\]
and $Q_p$ is the unique positive radial solution of the equation
\[
 -\Delta Q+Q=|Q|^{p-2}Q.
\]
\end{lemma}

\begin{lemma}[Gagliardo-Nirenberg inequality of Choquard type\cite{Mza}] \label{123}
Let $N\geq3$, $\alpha\in(0,N)$ and $\frac{N+\alpha}{N}<p<\frac{N+\alpha}{N-2}$, we have
\begin{equation}\label{GN2}
 \int_{\mathbb{R}^N}(I_\alpha\ast|u|^p)|u|^pdx
 \leq C_{N,p}\Big(\int_{\mathbb{R}^N}|\nabla u|
 ^2dx\Big)^{\frac{Np-(N+\alpha)}{2}}
 \Big(\int_{\mathbb{R}^N}|u|^2dx\Big)^{\frac{N+\alpha-p(N-2)}{2}},
\end{equation}
where equality holds for $u=Q_p$, $C_{N,p}=\frac{p}{|Q_p|_2^{2p-2}}$ and $Q_p$ 
is a nontrivial solution of
\begin{equation}
 -\frac{N(p-2)+N-\alpha}{2}\Delta Q_p+\frac{N+\alpha-(N-2)p}{2}Q_p
 =(I_\alpha\ast|Q_p|^p)|Q_p|^{p-2}Q_p.
\end{equation}
\end{lemma}
By \cite[(3.3)]{banqun}, we have the following Lemma.

\begin{lemma}\label{qxz3}
If $\frac{N+\alpha}{N}<r_1,r_2<\frac{N+\alpha}{N-2}$, then
\[
\int_{\mathbb{R}^N}(I_\alpha\ast|u_1|^{r_1})|u_2|^{r_2}dx
\leq\Big(\int_{\mathbb{R}^N}(I_\alpha\ast|u_1|^{r_1})|u_1|^{r_1}dx\Big)^{1/2}
\Big(\int_{\mathbb{R}^N}(I_\alpha\ast|u_2|^{r_2})|u_2|^{r_2}dx\Big)^{1/2},
\]
where $I_\alpha$ is given by \eqref{gamma}.
\end{lemma}


\begin{lemma}[{\cite[Lemma A.2]{em}}] \label{Liuv}
If $p\in(0,\frac{N}{N-2}]$ when $N\geq3$ and $p\in(0,\infty)$ when $N=1,2$. 
Let $u\in L^p(\mathbb{R}^N)$ be a smooth nonnegative function
satisfying $-\Delta u\geq0$ in $\mathbb{R}^N$. Then $u\equiv0$.
\end{lemma}

We now consider the problem
\begin{equation}\label{sig}
\begin{gathered}
-(a+b\int_{\mathbb{R}^N}|\nabla u|^2dx)\Delta u
=\lambda u+\mu(I_\alpha\ast|u|^p)|u|^{p-2}u,\\
\int_{\mathbb{R}^N}|u|^2dx=d>0,
\end{gathered}
\end{equation}
where $a,b>0$, $\frac{N+\alpha}{N}<p<\frac{N+\alpha+2}{N}$ and $N=3,4$. 
The corresponding minimization problem of \eqref{sig} is
\begin{equation}\label{ehs}
\sigma_p^{\mu}(d):=\inf_{u\in S(d)} G_\mu(u),
\end{equation}
where $G_\mu(u)=\frac{a}{2}\| \nabla u\|_2^2+\frac{b}{4}\| \nabla u\|_2^4
-\frac{\mu}{2p} \int_{\mathbb{R}^N}(I_\alpha\ast|u|^p)|u|^pdx$.

\begin{lemma}\label{byg}
\begin{itemize}
\item[(i)]   For any $d>0$, we have $\sigma_p^{\mu}(d)<0$.
\item[(ii)] $\sigma_p^{\mu}(d)$ is continuous with respect to $d\geq0$.
\item[(iii)] For any $d\geq m\geq0$, then 
$\sigma_p^{\mu}(d)\leq\sigma_p^{\mu}(m)+\sigma_p^{\mu}(d-m)$.
\end{itemize}
\end{lemma}

\begin{proof}
Item (i) follows from \cite[Theorem 1.2]{ygg}. 
For item (ii), we assume that $d^n=d+o_n(1)$. By the definition of
$\sigma_p^{\mu}(d^n)$, for any $\varepsilon>0$, there exists $u^n\in
 S_r(d^n)$ such that
\begin{equation}\label{gmb5}
 G_\mu(u^n)\leq\sigma_p^{\mu}(d^n)+\varepsilon.
\end{equation}
Let $w^n:=\frac{u^n}{\| u^n \|_2d^{1/2}}$. It is
easy to check that $w^n\in S_r(d^n)$ and
\begin{equation}\label{gmb6}
 \sigma_p^{\mu}(d)\leq G_\mu(w^n)=G_\mu(u^n)+o_n(1).
\end{equation}
Combining \eqref{gmb5} and \eqref{gmb6}, we have
\[
 \sigma_p^{\mu}(d)\leq\sigma_p^{\mu}(d^n)+\varepsilon+o_n(1).
\]
Reversing the argument, we obtain similarly that
\[
\sigma_p^{\mu}(d^n)\leq\sigma_p^{\mu}(d)+\varepsilon+o_n(1).
\]
Thus, by the arbitrariness of $\varepsilon$, we have
$\sigma_p^{\mu}(d^n)=\sigma_p^{\mu}(d)+o_n(1)$. This proves (ii).

By the density of $C_0^{\infty}(\mathbb{R}^N)$ in $H^1(\mathbb{R}^N)$, 
for any $\varepsilon>0$, there exist $\tilde{\psi}$, $\bar{\psi}
 \in C_0^{\infty}(\mathbb{R}^N)\times C_0^{\infty}(\mathbb{R}^N)$ with 
$\| \tilde{\psi}\|_2^2=m$, $\| \bar{\psi}\|_2^2=d-m$,
$i=1,2$ such that
\begin{gather*}
G_\mu(\tilde{\psi})\leq\sigma_p^{\mu}(m)+\varepsilon,\\
G_\mu(\bar{\psi})\leq\sigma_p^{\mu}(d-m)+\varepsilon.
\end{gather*}
Since $G_\mu$ is invariant by translation, without loss of generality, we may assume 
that $\operatorname{supp}\tilde{\psi}\cap \operatorname{supp}\bar{\psi}=\emptyset$, 
and then
 $\| \tilde{\psi}+\bar{\psi}\|_2^2=\| \tilde{\psi}\|_2^2+\|\bar{\psi}\|_2^2=d$,
\[
 \sigma_p^{\mu}(d)\leq G_\mu(\tilde{\psi}+\bar{\psi})\leq\sigma_p^{\mu}(m)+
 \sigma_p^{\mu}(d-m)+2\varepsilon.
\]
Therefore,
\[
 \sigma_p^{\mu}(d)\leq\sigma_p^{\mu}(m)+\sigma_p^{\mu}(d-m).
\]
This proof is complete.
\end{proof}


\begin{lemma}[{[?, Lemma 2.5]linjie}] \label{ph2}
Let $\frac{N+\alpha}{N}<r_1,r_2<\frac{N+\alpha}{N-2}$, if
\[
(u_1^n,u_2^n)\rightharpoonup(u_1,u_2) \text{in} H^1(\mathbb{R}^N)\times H^1(\mathbb{R}^N),
\]
then, up to a subsequence,
\begin{equation}\label{Blm}
 \int_{\mathbb{R}^N}(I_\alpha\ast|u_1^n|^{r_1})|u_2^n|^{r_2}dx
 =\int_{\mathbb{R}^N}(I_\alpha\ast|u_1|^{r_1})|u_2|^{r_2}dx+o_n(1).
\end{equation}
\end{lemma}

\begin{lemma}\label{fud}
If {\rm (A1)} holds, then $\inf_{W(k)}F(u_1,u_2)<0$ for all $k>0$.
\end{lemma}

\begin{proof}
Let $u^{\tau}(x)=\tau^{\frac{N}{2}}u(\tau x)$. Then it is easy to check that
\[
(u_1,u_2)\in S_r(d_1)\times S_r(d_2), \quad
 (u_1^{\tau},u_2^{\tau})\in W(k)
\]
when $\tau$ is sufficiently small.
From
\begin{align*}
F(u_1^{\tau},u_2^{\tau})
&= \frac{a}{2}\tau^2\sum_{i=1}^2\| \nabla u_i\|_2^2+\frac{b}{4}\tau^4\sum_{i=1}^2\|
 \nabla u_i\|_2^4-\sum_{i=1}^2\frac{\mu_i}{2p_i}\tau^{N{p_i}-N-\alpha}
  \int_{\mathbb{R}^N}(I_\alpha\ast|u_i|^{p_i})|u_i|^{p_i}dx\\
&\quad -\Theta\tau^{\frac{Nr}{2}-N-\beta}\int_{\mathbb{R}^N}(I_\beta\ast|u_1|^{p_1})
 |u_2|^{p_2}dx,
\end{align*}
$\frac{N+\alpha}{N}<p_1,p_2<\frac{N+\alpha+2}{N}$ and 
$\frac{N+\beta}{N}<r_1,r_2<2_\beta^*$, we have $\Phi_{(u_1,u_2)}(\tau) <0$ for
 $\tau$ small enough.
\end{proof}

\begin{lemma}\label{dx5}
Assume that {\rm (A2)} holds. Then there exist $k_0=k_0(d_1,d_2)>0$, 
$\Theta_*=\Theta_*(d_1,d_2)>0$ such that for any $0<\Theta<\Theta_\star$,
\[
\inf_{W(2k_0)\setminus W(k_0)}F(u_1,u_2)>0.
\]
And there exists $\varepsilon_0>0$ small enough, such that
\[
\sigma(d_1,d_2)<\inf_{\overline{W(k_0)}\setminus A(k_0-\varepsilon_0)}F(u_1,u_2).
\]
\end{lemma}

\begin{proof}
For  $(u_1,u_2)\in S_r(d_1)\times S_r(d_2)$, 
let $k=\| \nabla u_1\|_2^2+\| \nabla u_2\|_2^2$. Then
by Lemma \ref{123} and Lemma \ref{qxz3}, we have
\begin{equation}\label{dmk7}
\begin{aligned}
&F(u_1,u_2) \\
&=\frac{a}{2}k+\frac{b}{4}\sum_{i=1}^2\| \nabla u_i\|_2^4-
\sum_{i=1}^2\frac{\mu_i}{2p_i}\int_{\mathbb{R}^N}
 (I_\alpha\ast|u_i|^{p_i})|u_i|^{p_i}dx-\Theta\int_{\mathbb{R}^N}
 (I_\beta\ast|u_1|^{p_1})|u_2|^{p_2}dx\\
&\geq \frac{b}{8}k^2-\sum_{i=1}^2B_i\| \nabla u_i\|_2^{Np_i-N-\alpha}
 -\Theta B_3k^{\frac{Nr-2(N+\beta)}{4}}\\
&\geq \frac{b}{8}k^2-\sum_{i=1}^2B_ik^{\frac{Np_i-N-\alpha}{2}}
 -\Theta B_3k^{\frac{Nr-2(N+\beta)}{4}}:=g(k),\\
\end{aligned}
\end{equation}
where $B_i=\frac{\mu_i}{2p_i}C_{N,p_i}d_i^{\frac{N+\alpha-p_i(N-2)}{2}}$ ($i=1,2$), 
and $B_3=C_{N,r_1}C_{N,r_2}d_1^{\frac{N+\beta-r_1(N-2)}{4}}
d_2^{\frac{N+\beta-r_2(N-2)}{4}}$.
By (A2), we have
\[
0<\frac{Np_i-N-\alpha}{2}<1 \quad \text{and}\quad
  \frac{Nr-2(N+\beta)}{4}>2, \quad i=1,2.
\]
We fix $k=k_0$ sufficiently large such that
\[
\sum_{i=1}^2B_ik_0^{\frac{Np_i-N-\alpha-4}{2}}\leq\frac{b}{24},
\]
and fix $\Theta=\Theta_*$ sufficiently small such that
\[
\Theta_*B_3(2k_0)^{\frac{Nr-2(N+\beta)-8}{4}}\leq\frac{b}{24}.
\]
Therefore, for each $0<\Theta\leq\Theta_*$ and
$(u_1,u_2)\in W(2k_0)\setminus W(k_0)$,
we obtain
\begin{align*}
F(u_1,u_2)
&\geq \frac{b}{8}k^2-\sum_{i=1}^2B_ik^{\frac{Np_i-N-\alpha}{2}}
 -\Theta B_3k^{\frac{Nr-2(N+\beta)}{4}}\\
&= k^2\Big(\frac{b}{8}-\sum_{i=1}^2B_ik^{\frac{Np_i-N-\alpha-4}{2}}
 -\Theta B_3k^{\frac{Nr-2(N+\beta)-8}{4}}\Big)\\
&\geq bk_0^2\Big(\frac{1}{8}-\frac{1}{24}-\frac{1}{24}\Big)
=\frac{b}{24}k_0^2.
\end{align*}

Next, by continuity of $g(k)$ and $g(k_0)>0$, there exists $\varepsilon_0>0$ 
sufficiently small such that $g(k)\geq0$ when $k\in[k_0-\varepsilon,k_0]$. Hence,
\[
F(u_1,u_2)\geq g(k)\geq0>\sigma(d_1,d_2)
\]
for any $(u_1,u_2)\in\overline{W(k_0)}\setminus W(k_0-\varepsilon_0)$.
This proof is complete.
\end{proof}

\begin{lemma}\label{zjngz}
Let $\frac{N+\alpha}{N}<p_1,p_2<2_\alpha^*$ and $\frac{N+\beta}{N}<r_1,r_2<2_\beta^*$. 
If $(u_1,u_2)\neq(0,0)$ is a solution of
 \eqref{yy} for some $(\lambda_1,\lambda_2)\in\mathbb{R}^2$, then $\lambda_1<0$ 
 or $\lambda_2<0$. Furthermore, if $u_1>0$ and $u_2\geq0$, then
$\lambda_1<0$; if $u_1\geq0$ and $u_2>0$, then $\lambda_2<0$.
\end{lemma}

\begin{proof}
Testing \eqref{yy} by $(u_1,u_2)$ and integrating in $\mathbb{R}^N$, one has
\begin{align*}
&\lambda_1\| u_1\|_2^2+\lambda_2\| u_2\|_2^2 \\
&=a\sum_{i=1}^2\| \nabla u_i\|_2^2+
b\sum_{i=1}^2\| \nabla u_i\|_2^4-\sum_{i=1}^2\mu_i\int_{\mathbb{R}^N}(I_\alpha\ast|u_i|^{p_i})|u_i|^{p_i}dx-
\Theta r\int_{\mathbb{R}^N}(I_\beta\ast|u_1|^{r_1})|u_2|^{r_2}dx.
\end{align*}
Combining $P(u_1,u_2)=0$, $2<p_1,p_2<2^*$ and $\frac{N+\beta}{N}<r_1,r_2<2_\beta^*$,
we have
 \begin{align*}
 \lambda_1\| u_1\|_2^2+\lambda_2\| u_2\|_2^2
&= \sum_{i=1}^2\mu_i(Np_i-N-\alpha-1)
 \int_{\mathbb{R}^N}(I_\alpha\ast|u_i|^{p_i})|u_i|^{p_i}dx\\
&\quad +\Theta\left(\frac{Nr-2(N+\beta)}{4}-r\right)\int_{\mathbb{R}^N}(I_\beta\ast|u_1|^{p_1})|u_2|^{p_2}dx<0.
 \end{align*}
Thus, at least one of $\lambda_1$ and $\lambda_2$ is negative.

Then, we argue by contradiction, suppose $\lambda_1\geq0$. 
In view $u_1>0$, we have
\[
-\Big(a+b\int_{\mathbb{R}^N}|\nabla u_1|^2dx\Big)\Delta u_1
=\lambda_1u_1+\mu_1(I_\alpha\ast|u_1|^{p_1})|u_1|^{p_1-2}u_1
 +\Theta r_1(I_\beta\ast|u_2|^{p_2})|u_1|^{p_1-2}u_1\geq0.
\]
By Lemma \ref{Liuv}, we obtain $u_1=0$, which  contradicts $u_1\neq0$.
Therefore, $\lambda_1<0$. The other case can be proved
in the same way.
\end{proof}

\begin{lemma}\label{zyy4}
Assume that $2<p_1,p_2<\frac{2N}{N-2}$ and 
$\frac{N+\beta}{N}<r_1,r_2<\frac{N+\beta}{N-2}$. For any bounded Palais-Smale sequence
$\{(u_1^n,u_2^n)\}$ for $F$ on $S_r(d_1)\times S_r(d_2)$, then there exist 
$(u_1,u_2)\in H_r^1(\mathbb{R}^N)\times H_r^1(\mathbb{R}^N)$ and a sequence
$\{(\lambda_1^n,\lambda_2^n)\}\subset\mathbb{R}^2$, such that up to a subsequence
\begin{itemize}
\item[(a)] $(u_1^n,u_2^n)\rightharpoonup (u_1,u_2)$ in 
 $H_r^1(\mathbb{R}^N)\times H_r^1(\mathbb{R}^N)$, 
 $(u_1^n,u_2^n)\to  (u_1,u_2)$ in $L^p
(\mathbb{R}^N)\times L^p(\mathbb{R}^N)$ for $p\in(2,\frac{2N}{N-2})$.

\item[(b)] $(\lambda_1^n,\lambda_2^n) \to  (\lambda_1,\lambda_2)$ in
 $(\mathbb{R}^2)$.
 
\item[(c)] $F'(u_1^n,u_2^n)-\lambda_1^n(u_1^n,0)-\lambda_2^n(0,u_2^n)\to 0$
 in $H_r^{-1}(\mathbb{R}^N)\times H_r^{-1}(\mathbb{R}^N)$.
 
\item[(d)] $(u_1^n,u_2^n)$ is a solution to the system \eqref{yy} for some 
 $\lambda_1,\lambda_2\leq0$ if $(u_1,u_2)$ satisfies the additional property
$P(u_1^n,u_2^n)\to 0$, where $(\lambda_1,\lambda_2)$ is given by (b).
\end{itemize}
\end{lemma}

\begin{proof}
Obviously, item (a) is true 
Since $\{(u_1^n,u_2^n)\}\subset H_r^1(\mathbb{R}^N)\times H_r^1(\mathbb{R}^N)$ 
is bounded, by \cite{Arm2},
we have $(F|_{S_r(d_1)\times S_r(d_2)})'(u_1^n,u_2^n)\to 0$ in 
$H_r^{-1}(\mathbb{R}^N)\times H_r^{-1}(\mathbb{R}^N)$ is
equivalent to
\[
F'(u_1^n,u_2^n)-\frac{1}{\| u_1^n\|_2^2}\langle F'(u_1^n,u_2^n),
(u_1^n,0)\rangle(u_1^n,0)-
\frac{1}{\| u_2^n\|_2^2}\langle F'(u_1^n,u_2^n),(0,u_2^n)\rangle(0,u_2^n)\to 0
\]
in $H_r^{-1}(\mathbb{R}^N)\times H_r^{-1}(\mathbb{R}^N)$. Thus, we have
\[
 F'(u_1^n,u_2^n)-\lambda_1^n(u_1^n,0)-\lambda_2^n(0,u_2^n)\to 0
 \quad \text{in } H_r^{-1}(\mathbb{R}^N)\times H_r^{-1}(\mathbb{R}^N)
\]
with
\begin{equation}\label{yjj1}
\begin{aligned}
\lambda_1^n&=\frac{1}{\| u_1^n\|_2^2}\Big(a\| \nabla u_1^n\|_2^2+b\| \nabla u_1^n\|_2^4-
\mu_1\int_{\mathbb{R}^N}(I_\alpha\ast|u_1^n|^{p_1})|u_1^n|^{p_1}dx \\
&\quad -\Theta r_1\int_{\mathbb{R}^N}(I_\beta\ast|u_1^n|^{p_1})|u_2^n|^{p_2}dx\Big)
-o_n(1),
\end{aligned}
\end{equation}
and
\begin{equation}\label{yjj2}
\begin{aligned}
 \lambda_2^n
&=\frac{1}{\| u_2^n\|_2^2}\Big(a\| \nabla u_2^n\|_2^2+b\| \nabla u_2^n\|_2^4-
 \mu_2\int_{\mathbb{R}^N}(I_\alpha\ast|u_2^n|^{p_2})|u_2^n|^{p_2}dx \\
&\quad -\Theta r_2\int_{\mathbb{R}^N}(I_\beta\ast|u_1^n|^{p_1})|u_2^n|^{p_2}dx\Big)-o_n(1).
\end{aligned}
 \end{equation}
This proves (c).
By the boundedness of $u_1^n,u_2^n$ in $H_r^1(\mathbb{R}^N)$, $L^p(\mathbb{R}^N)$
for $p\in(2,2^*)$, Lemma \ref{123} and Lemma \ref{qxz3},
we obtain $\{\lambda_1^n\}$, $\{\lambda_2^n\}$ are bounded. This proves (b).
Combining (b) and (c), it is now standard to deduce (d).
\end{proof}

\begin{lemma}\label{yuk9}
Under the conditions of Lemma \ref{zyy4},  $u_1^n\to  u_1$ in $H_r^1(\mathbb{R}^N)$
if $\lambda_1<0$. Similarly,
$u_2^n\to  u_2$ in $H_r^1(\mathbb{R}^N)$ if $\lambda_2<0$.
\end{lemma}

\begin{proof}
In view of Lemmas \ref{ph2} and  \ref{zyy4}, we obtain
\begin{gather}\label{vbp1}
\int_{\mathbb{R}^N}(I_\alpha\ast|u_i^n|^{p_i})|u_i^n|^{p_i}dx
\to \int_{\mathbb{R}^N}(I_\alpha\ast|u_i|^{p_i})|u_i|^{p_i}dx 
(\frac{N+\alpha}{N}<p_i<2_\alpha^*,\quad  i=1,2), \\
\label{vbp2}
\int_{\mathbb{R}^N}(I_\beta\ast|u_1^n|^{r_1})|u_2^n|^{r_2}dx
\to \int_{\mathbb{R}^N}(I_\beta\ast|u_1|^{r_1})|u_2|^{r_2}dx (\frac{N+\beta}{N}
<r_1,r_2<2_\beta^*), \\
\langle F'(u_1^n,u_2^n)-\lambda_1^n(u_1^n,0),(u_1^n,0)\rangle\to 0
= \langle F'(u_1,u_2)-\lambda_1^n(u_1,0),(u_1,0)\rangle.
\end{gather}
Therefore,
\[
a\| \nabla u_1^n\|_2^2+b\| \nabla u_1^n\|_2^4-\lambda_1^n\| u_1^n\|_2^2\to
a\| \nabla u_1\|_2^2+b\| \nabla u_1\|_2^4-\lambda_1\| u_1\|_2^2.
\]
Since
\begin{equation}\label{vbp3}
 \| u_1\|_2^2\leq \varliminf_{n\to +\infty} \| u_1^n\|_2^2, \quad
 \| \nabla u_1\|_2^2\leq \varliminf_{n\to +\infty} \| \nabla u_1^n\|_2^2,
\end{equation}
it follows that
\[
 \| \nabla u_1^n\|_2^2\to \| \nabla u_1\|_2^2, \quad
 \| u_1^n\|_2^2\to \| u_1\|_2^2.
\]
This proof is complete.
\end{proof}

\section{Proof of Theorem \ref{thm1.1}} 

\begin{lemma}\label{xj5}
Assume that {\rm (A1)} holds. Then $F$ is bounded form below and coercive on 
$S_{r}(d_1)\times S_{r}(d_2)$. Furthermore, there exists a bounded
Palais-Smale sequence $\{(u_1^n,u_2^n)\}\subset S_{r}(d_1)\times S_{r}(d_2)$, 
which satisfies $(u_1^n)^-\to 0$ and
$(u_2^n)^-\to 0$ in $H_r^1(\mathbb{R}^N)\times H_r^1(\mathbb{R}^N)$.
\end{lemma}

\begin{proof}
Since $\frac{N+\alpha}{N}<p_1,p_2<\frac{N+\alpha+2}{N}$, 
$2\cdot\frac{N+\beta}{N}<r_1+r_2<2\cdot\frac{N+\beta+4}{N}$, we have $0<Np_i-N-\alpha<2$ 
and $\frac{Nr-2(N+\beta)}{4}<2$.
From \eqref{dmk7}, we know that $F$ is bounded form below and coercive on 
$S_{r}(d_1)\times S_{r}(d_2)$.

Next, let $\{(u_1^n,u_2^n)\}\subset S_{r}(d_1)\times S_{r}(d_2)$ be a minimizing 
sequence for $F$. By the coerciveness of $F$, the sequence is
bounded. Since the functional $F$ is even, we may assume that $w_1^n\geq0$ and 
$w_2^n\geq0$. It is easy to check that $F$ is a $C^1-$manifold in
$H_r^1(\mathbb{R}^N)\times H_r^1(\mathbb{R}^N)$, then by Ekeland's variational 
principle, there exists a minimizing sequence
$\{(u_1^n,u_2^n)\}\subset S_{r}(d_1)\times S_{r}(d_2)$, which is the Palais-Smale 
sequence for $F$ restricted to $S_{r}(d_1)\times S_{r}(d_2)$ and satisfies
$\| (u_1^n,u_2^n)-(w_1^n,w_2^n)\|_{H_r^1(\mathbb{R}^N)\times H_r^1(\mathbb{R}^N)}\to 0$ 
as $n\to +\infty$.
From $w_1^n\geq0$ and $w_2^n\geq0$, we have $(u_1^n)^-\to 0$ and $(u_2^n)^-\to 0$ in 
$H_r^1(\mathbb{R}^N)\times H_r^1(\mathbb{R}^N)$.
\end{proof}


\begin{proof}[Proof of Theorem \ref{thm1.1}] 
By Lemma \ref{xj5}, there exists a bounded Palais-Smale sequence 
$\{(u_1^n,u_2^n)\}\subset S_{r}(d_1)\times S_{r}(d_2)$
and $(u_1,u_2)\in H_r^1(\mathbb{R}^N)\times H_r^1(\mathbb{R}^N)$ with $u_1\geq0$ 
and $u_2\geq0$, such that $\lim_{n\to +\infty}
(u_1^n,u_2^n)=\sigma(d_1,d_2)$, $(u_1^n,u_2^n)\rightharpoonup(u_1,u_2)$ in 
$H_r^1(\mathbb{R}^N)\times H_r^1(\mathbb{R}^N)$. Hence, it suffices to
prove that $(u_1^n,u_2^n)\to (u_1,u_2)$ in 
$H_r^1(\mathbb{R}^N)\times H_r^1(\mathbb{R}^N)$. Indeed, if this holds, then we have
$\{(u_1^n,u_2^n)\}\subset S_{r}(d_1)\times S_{r}(d_2)$ and 
$F(u_1,u_2)=\sigma(d_1,d_2)$. Furthermore, by the strong maximum principle,
we have $u_1>0$ and $u_2>0$. From $\Theta>0$, obviously, we have 
$\sigma(d_1,d_2)\leq\sigma_{p_1}^{\mu_1}(d_1)+\sigma_{p_2}^{\mu_2}(d_2)$. 
By Lemma \ref{fud}, we know
$\sigma(d_1,d_2)<0$. We divide four cases to have that $u_1>0$ and $u_2>0$.
\smallskip 

\noindent\textbf{Case I:} $u_1=0$ and $u_2=0$.
Obviously, we can obtain that $\lim_{n\to +\infty}F(u_1^n,u_2^n)=0$, which is 
 contradicts $\sigma(d_1,d_2)<0$.
\smallskip

\noindent\textbf{Case II:} $u_1=0$ and $u_2\neq0$.
Combining \eqref{vbp1} and \eqref{vbp3}, we have
\begin{equation}\label{kgf}
\varliminf_{n\to +\infty}F(u_1^n,u_2^n)
\geq\frac{a}{2}\| \nabla u_2\|_2^2+\frac{b}{4}\|
 \nabla u_2\|_2^4-\frac{\mu_2}{2p_2}
 \int_{\mathbb{R}^N}(I_\alpha\ast|u_2|^{p_2})|u_2|^{p_2}dx
\geq \sigma_{p_2}^{\mu_2}(d_3),
\end{equation}
where $d_3:=\| u_2\|_2^2\leq d_2$. By (i) and (iii) of Lemma \ref{byg}, 
we have $\sigma_{p_2}^{\mu_2}(d_3)\geq\sigma_{p_2}^{\mu_2}(d_2)$.
Thus, by \eqref{kgf}, we have $\sigma(d_1,d_2)\geq \sigma_{p_2}^{\mu_2}(d_3)$. 
On the other hand, using Lemma \ref{byg} again, we have $\sigma_{p_1}^{\mu_1}(d_1)<0$. 
In view of $\sigma(d_1,d_2)\leq\sigma_{p_1}^{\mu_1}(d_1)+\sigma_{p_2}^{\mu_2}(d_2)$, 
we have $\sigma(d_1,d_2)<\sigma_{p_2}^{\mu_2}(d_2)$.
This is a contradiction.
\smallskip

\noindent\textbf{Case III:} $u_1\neq0$ and $u_2=0$.
By the same proof as in Case II, we can  obtain a contradiction.
\smallskip

\noindent\textbf{Case IV:} $u_1\neq0$ and $u_2\neq0$.
From Lemma \ref{zjngz}, we have $\lambda_1<0$ and $\lambda_2<0$. Thus, by 
Lemma \ref{yuk9}, we obtain $u_1^n\to  u_1$ and
$u_2^n\to  u_2$ in $H_r^1(\mathbb{R}^N)$. This proof is complete.
\end{proof}

\begin{thebibliography}{99}
\bibitem[1]{Kf1} A. Arosio, S. Panizzi;
\emph{On the well-posedness of the Kirchhoff string}. 
Trans. Am. Math. Soc., 348 (1996), 305-330.
\bibitem[2]{Huij4} T. Bartsch, H. W. Li, W. M. Zou;
\emph{Existence and asymptotic behavior of normalized ground states for Sobolev
critical Schr\"{o}dinger systems}, Calc. Var. Partial Differential Equations, 62 (2023) 9.
\bibitem[3]{Huij1} T. Bartsch, L. Jeanjean, N. Soave;
\emph{Normalized solutions for a system of coupled cubic Schr\"{o}dinger equations
on $\mathbb{R}^3$}, J. Math. Pures Appl., 106 (2016) 583-614.
\bibitem[4]{Huij5} T. Bartsch, N. Soave;
\emph{A nature constraint approach to normalized solutions of nonlinear 
Schr\"{o}dinger equations and systems}, J. Funct. Anal., 272 (2017) 4998-5037.
\bibitem[5]{Huij2} T. Bartsch, N. Soave;
\emph{Multiple normalized solutions for a competing system of Schr\"{o}dinger equations},
 Calc. Var. Partial Differential Equations 58 (2019) 22.
\bibitem[6]{Huij3} T. Bartsch, X. X. Zhong, W. M. Zou;
\emph{Normalized solutions for a coupled Schr\"{o}dinger system}, Math. Ann.,
380 (2021) 1713-1740.
\bibitem[7]{Arm2} H. Berestycki, P. L. Lions;
\emph{Nonlinear scalar field equations II},
 Arch. Ration. Mech. Anal., 82(1983) 347-376.
\bibitem[8]{kf2} X. Cao, J. Xu, J. Wang;
\emph{The existence of solutions with prescribed $L^2$-norm for Kirchhoff type system}. 
J. Math. Phys., 58 (2017) 041502.
\bibitem[9]{BZd27} I. E. Dzyaloshinskii, E. M. Lifshitz, Lev P. Pitaevskii;
\emph{General theory of van der Waals' forces}. Sov. Phys. Usp., 4 (1961) 153.
\bibitem[10]{JFx} Q. P. Geng, Y. Y. Tu, J. Wang;
\emph{Existence and multiplicity of the positive normalized solutions to the
coupled Hartree-Fock type nonlocal elliptic system}, 
J. Fixed Point Theory Appl., 24 (2022) 83.
\bibitem[11]{banqun} M. Ghimenti, J. Van Schaftingen;
\emph{Nodal solutions for the Choquard equation}, J. Funct. Anal. 271
(2016) 107-135.
\bibitem[13]{Huij6} T. X. Gou, L. Jenajean;
\emph{Multiple positive normalized solutions for nonlinear Schr\"{o}dinger systems},
Nonlinearity, 31 (2018) 2319-2345.
\bibitem[14]{guo1} Z. Y. Guo, L. J. Zhao;
\emph{Existence of ground states for fractional Choquard-Kirchhoff
equations with magnetic fields and critical exponents}, 
Period. Math. Hungar. 87 (2023) 468-483.
\bibitem[15]{guo2} Z. Y. Guo, T. Q. Zhang;
\emph{Normalized solutions of fractional Kirchhoff equations in the
defocusing case}, Electronic J. Differ. Equ., 2025 (2025) No. 75, 1-16.
\bibitem[16]{Kf2} J. Q. Hu, A. M. Mao;
\emph{Normalized solutions to nonlocal Schr\"{o}dinger systems with $L^2$-subcritical
 and supercritical nonlinearities}, J. Fixed Point Theory Appl., 25 (2023), 77.
\bibitem[17]{Kf5} Q. H. He, Z. Y. Lv, Z. W. Tang;
\emph{The existence of normalized solutions to the Kirchhoff equation with potential 
and Sobolev critical nonlinearities}, J. Geom. Anal., 33 (2023), 236.
\bibitem[18]{Kf6} Q. H. He, Z. Y. Lv, Y. M. Zhang, X. X. Zhong;
\emph{Existence and blow up behavior of positive normalized solution to the Kirchhoff
 equation with general nonlinearities:  Mass super-critical case},
  J. Differ. Equ., 356 (2023), 375-406.
\bibitem[19]{em} N. Ikoma;
\emph{Compactness of minimizing sequences in nonlinear Schr\"{o}dinger systems under
multiconstraint conditions}, Adv. Nonlinear Stud., 14 (2014), 115-136.
\bibitem[21]{Kf4} G. Kirchhoff;
\emph{Mechanik}. Teubner, Leipzig, 1883.
\bibitem[22]{lieb} E. H. Lieb, M. Loss;
\emph{Analysis}, Gradute Studies in Mathematics, AMS, Providence, Rhode island,
2001.
\bibitem[23]{Li4} J. Lions;
\emph{On some questions in boundary value problems of mathematical physics}. 
North-Holland Math. Stud., 30 (1978), 284-346.
\bibitem[24]{ygg} Z. Liu;
\emph{Multiple normalized solutions for Choquard equations involving Kirchhoff type 
perturbation}, Topol. Methods Nonlinear Anal., 54 (2019) 297-319.
\bibitem[26]{Mza} V. Moroz, J. Van Schaftingen;
\emph{Existence of groundstates for a class of nonlinear Choquard
equations}, Trans. Amer. Maths. Soc., 367 (2015), 6557-6579.
\bibitem[27]{cfrv} S. G. Porsev, A. Derevianko;
\emph{High-accuracy calculations of dipole, quadrupole, and octupole electric dynamic 
polarizabilities and van der Waals coefficients $C_6$, $C_8$, and $C_{10}$ 
for alkaline-earth dimers}, J. Exp. Theor. Phys. 102 (2006), 195-205.
\bibitem[28]{GN23} M. I. Weinstein;
\emph{Nonlinear Schr\"{o}dinger equations and sharp interpolation estimates}.
Comm. Math. Phys. 87 (1983), 567-576.
\bibitem[29]{Kf3} Z. Yang;
\emph{Normalized ground state solutions for Kirchhoff type systems}. 
J. Math. Phys., 62 (2021), 031504.
\bibitem[30]{VWp3} Y. Zheng, A. Narayanaswamy;
\emph{Lifshitz theory of van der Waals pressure in dissipative media}. 
Phys. Rev. A, 83 (2011), 042504.
\bibitem[31]{linjie} H. Zhang, J. J. Zhang, X. X. Zhong;
\emph{Normalized solutions for critical Choquard systems}, arXiv:2307.01483v2
\end{thebibliography}
\end{document}
